\section{Method of Moments and Maximal Likelihood Estimation}

\subsection{MoM}

\begin{definition}[Method of moments]
Let $X_1,\ldots,X_n$ be i.i.d.\ real observations from $P_\theta$,
where $\theta\in\Theta\subseteq\mathbb{R}^k$, and assume
$\mathbb{E}_\theta|X_1|^k<\infty$.
The \emph{population moments} and \emph{empirical moments} are
\[
\mu_j(\theta)=\mathbb{E}_\theta X_1^j,
\qquad
\widehat\mu_j=\frac{1}{n}\sum_{i=1}^n X_i^j,
\qquad 1\leq j\leq k.
\]

Write
\[
h(\theta)=\bigl(\mu_1(\theta),\ldots,\mu_k(\theta)\bigr),
\qquad
\widehat\mu=(\widehat\mu_1,\ldots,\widehat\mu_k).
\]
The \emph{method of moments} (MoM) solves
\[
h(\widehat\theta)=\widehat\mu
\]
for $\widehat\theta\in\Theta$. If $h:\Theta\to h(\Theta)$ is
one-to-one and the observed $\widehat\mu$ lies in $h(\Theta)$, then
\[
\widehat\theta=h^{-1}(\widehat\mu).
\]
As with every estimator, the inverse or the chosen solution must
be measurable as a function of the data.
\end{definition}

\subsection{MLE}

\begin{definition}[Maximum likelihood estimator]
Suppose $Y$ has a joint density or pmf $f_\theta$ relative to a
fixed measure that does not depend on $\theta$.
The \emph{likelihood} is
\[
L(\theta;y)=f_\theta(y),
\]
viewed as a function of $\theta$ with $y$ fixed.
A \emph{maximum likelihood estimator} (MLE) is a measurable choice
\[
\widehat\theta(y)
\in\operatorname*{arg\,max}_{\theta\in\Theta}L(\theta;y).
\]

The maximizing set may be empty or contain several points.
The likelihood is not a probability distribution over parameters.
Because the logarithm is increasing, we can instead maximize
the \emph{log-likelihood}
\[
\ell(\theta;y)=\log L(\theta;y),
\]
with the convention $\log 0=-\infty$.
\end{definition}
